\subsection{Environmental Selection Effects}
  \label{subsec:selection_effects}

  Local measurements of the Hubble constant are not uniformly distributed across cosmic environments.
  Distance ladder calibrators and standard candles are preferentially observed
  in specific regions of the large-scale structure~\cite{Riess2022SH0ES}.

  This selection effect plays a central role in the present framework.
  If observational strategies preferentially sample void-dominated or underdense regions,
  the inferred expansion rate is expected to be systematically enhanced relative to the global value.

  Importantly, this effect is not introduced as an ad hoc assumption.
  It constitutes a direct and testable prediction.
  Future analyses correlating inferred values of $H_0$ with void catalogs and density
  reconstructions can directly assess the magnitude of the effect, as already
  explored in the context of local underdensity scenarios~\cite{Shanks2019Void,Kennworthy2019VoidCritique}.
